\chapter{Fluxys}

\section{Overview of Fluxys: A Purpose-Driven Energy Infrastructure Company}

\subsection{Company Profile}

Fluxys is a Belgian-headquartered multi-molecule infrastructure company operating across the full gas value chain. With approximately 1,400 employees, the group maintains a strong European presence complemented by associated companies in South America and the Middle East. Its corporate mission is encapsulated in the tagline \textit{``shaping together a bright energy future''}, reflecting its positioning as a purpose-driven actor committed to accelerating the energy transition.

The ownership structure of Fluxys is characterised by a public-private partnership model designed to support long-term investment horizons. The principal shareholders of Fluxys Belgium include Publigas (77.38\%), Energy Infrastructure Partners -- EIP (15.22\%), the Federal Holding and Investment Company (3.44\%), AG Insurance (1.97\%), Ethias (1.31\%), and personnel and management (0.67\%). The Belgian State retains a golden share. Fluxys Belgium is also partly listed on Euronext (10\% of shares). Fluxys holds 100\% of both Fluxys Europe and Fluxys International.

Fluxys operates across the full gas value chain in Belgium, covering terminalling (regasification and liquefaction at the Zeebrugge LNG terminal), transmission, and underground storage (at Loenhout). Its Belgian network plays a structurally important role in the security of supply of Northwest Europe, acting as a transit hub connecting Norway, LNG imports, the United Kingdom, France, the Netherlands, Germany, and Luxembourg.

\subsection{Belgian Gas Network: Key Figures (2025)}

In 2025, total gas imports via pipeline and LNG into Belgium reached 480~TWh (up from 407~TWh in 2024), of which 147~TWh came from LNG terminals. Transit to surrounding countries amounted to 331~TWh. Domestic consumption stood at 151~TWh. The Loenhout underground storage facility has a storage capacity of 7.6~TWh, a send-out capacity of 7.25~GWh/h, and an injection capacity of 3.75~GWh/h; it reached 100\% filling level at the start of August 2025. The Zeebrugge LNG terminal has a regasification capacity of 197~TWh/year (17~bcm/year), LNG storage of 566,000~m\textsuperscript{3}, and sent 146.8~TWh into the network in 2025, up from 81.5~TWh in 2024.

\section{The European Energy Context and Regulatory Framework}

\subsection{Geopolitical Disruptions and Supply Shifts}

The Russian invasion of Ukraine in February 2022 triggered a dramatic restructuring of European gas supply. North Stream flows were halved in July 2022, reduced to 20\% of capacity, and halted definitively in August 2022. The Nord Stream sabotage occurred in September 2022. As of January 2025, Russian gas flows through Ukraine have fully ceased. In response, Europe has massively scaled up LNG imports: LNG terminal throughput across Europe grew from 915~TWh in 2021 to 1,614~TWh in 2025. Norwegian pipeline flows have remained relatively stable at around 1,233~TWh. These shifts have made Belgium's Zeebrugge terminal a critical entry point for European gas supply diversification.

Gas prices (TTF) spiked sharply during the 2022 crisis, reaching above 300~€/MWh, before normalising. More recent price movements have been driven by events including the end of Russian Ukrainian transit flows (January 2025), the second Trump presidency and associated geopolitical uncertainty, the Israel-Iran war (June 2025), a cold snap in Europe (January 2026), and Winter Storm Fern in Texas (January 2026).

\subsection{Regulatory Framework}

Gas (and hydrogen) transmission in Belgium is regulated, with Belgian legislation deriving primarily from EU law. The regulatory framework has evolved through successive EU energy packages:

\begin{itemize}[noitemsep]
    \item \textbf{1st and 2nd Energy Packages (1996--2004):} Liberalisation of gas markets, creation of national regulatory agencies.
    \item \textbf{3rd Energy Package (2009):} Third-party access, ownership unbundling, creation of ACER and ENTSOG.
    \item \textbf{Gas Network Codes (2013--2018):} Capacity allocation, interoperability, balancing, tariffs.
    \item \textbf{RePowerEU and related regulation (2022):} Phase-out of Russian gas dependency, mandatory storage filling targets, joint gas purchase platform.
    \item \textbf{H2 and Decarbonisation Package (July 2024):} Updated EU gas market rules and a new regulatory framework for hydrogen infrastructure, including tariffs, cost mutualization, inter-temporal cost allocation, and creation of ENNOH.
\end{itemize}

Key principles of EU gas regulation include ownership unbundling, regulated tariffs, third-party access (non-discrimination), transparency obligations (REMIT, ENTSOG Transparency Platform), network codes developed by ENTSOG and ACER, and a Ten-Year Network Development Plan (TYNDP).

\subsection{Tariff Computation}

Allowed revenues for regulated activities are computed as the sum of operating costs (OPEX), depreciation, and a regulated margin based on the Regulatory Asset Base (RAB). The RAB of Fluxys Belgium stands at approximately €2~billion. The regulated return is calculated via a WACC formula (currently 4.80\% gross), which combines:
\begin{itemize}[noitemsep]
    \item ROE on the first part of equity (up to 40\% of RAB), based on the CAPM model;
    \item ROEE on the second part of equity (above 40\% of RAB), remunerated slightly above the risk-free rate;
    \item Cost of debt on 45\% of RAB (pass-through mechanism).
\end{itemize}

Total allowed revenues amount to approximately €420~million, decomposed into €100~million return, €230~million operating costs, and €90~million depreciation. A regulatory account adjustment of $-$130~million is applied in the current tariff period. Unit tariffs per service are derived by dividing allowed revenues by reference quantities (Qref) for each activity (transmission, storage, LNG terminal). The tariff period is four years.

\section{The Energy Transition: Fluxys's Strategic Role}

\subsection{The Energy Trilemma}

As an essential infrastructure provider, Fluxys frames its energy transition strategy around balancing three imperatives: \textbf{sustainability} (reducing carbon emissions), \textbf{affordability} (keeping energy costs manageable), and \textbf{security of supply} (ensuring reliable access to energy). These three dimensions constitute the classical ``energy trilemma'' and guide Fluxys's infrastructure investment decisions.

\subsection{The EU Net-Zero Pathway and Belgium's Energy Mix}

The EU Green Deal targets net-zero greenhouse gas emissions by 2050, starting from 3.75~Gt CO\textsubscript{2} in 1990. Achieving this requires a combination of solutions: energy efficiency, green electricity generation, scaling up the hydrogen economy, biomethane and synthetic methane, and carbon capture and utilisation/storage (CCUS).

Belgium's 2024 final energy consumption stood at 435~TWh, dominated by oil (48\%), natural gas (25\%), and electricity (17\%). The Belgian electricity mix (81~TWh) was composed of nuclear (42\%), renewables and biofuels (33\%), natural gas (18\%), and other sources (8\%). Belgium's gas system has a peak capacity of 57~GW, far exceeding the electricity system's 14~GW, making gas infrastructure an indispensable source of seasonal and daily flexibility, especially as intermittent renewables expand.

\subsection{Ten-Year Network Development Plan and Investments}

Fluxys Belgium's 2025--2034 investment plan amounts to an indicative €8.3~billion. Over 90\% of this plan is directed towards infrastructure for low-carbon molecules and reducing the company's own greenhouse gas emissions. The main investment categories are: hydrogen infrastructure (the largest share), CO\textsubscript{2} infrastructure, facilitating green gas integration into the natural gas network, reducing Fluxys's own GHG emissions, and other investments.

Hydrogen and CO\textsubscript{2} transmission in Belgium are managed through three separate Fluxys Belgium affiliates: \textit{Fluxys Hydrogen} (100\% owned, nominated as the Belgian Hydrogen Network Operator for 20 years), \textit{Fluxys c-grid} (77.5\% owned, CO\textsubscript{2} network operator in Wallonia and Flanders), and \textit{Fluxys c-grid Antwerp} (70\% owned, CO\textsubscript{2} network operator application submitted for the Antwerp port cluster).

\subsection{Biomethane}

Biomethane is highlighted as a ``no-brainer'' near-term decarbonisation lever. Produced from organic waste, livestock effluents, and agriculture-related feedstocks through anaerobic digestion and purification, biomethane is chemically identical to natural gas and can be injected directly into the existing grid. The EU27+UK biomethane supply potential is estimated at 1,110~TWh/year by 2050, covering applications in transport (566~TWh), power (300~TWh), buildings (178~TWh), and industry (66~TWh).

\section{Hydrogen Infrastructure}

\subsection{Hydrogen Colours and Production Pathways}

Hydrogen is produced through several routes with varying carbon intensities:
\begin{itemize}[noitemsep]
    \item \textbf{Grey hydrogen:} Steam methane reforming (SMR) without carbon capture.
    \item \textbf{Blue hydrogen:} SMR or autothermal reforming with carbon capture -- classified as Low Carbon H2.
    \item \textbf{Green hydrogen (RFNBO):} Electrolysis powered by renewable electricity (wind or solar).
    \item \textbf{Pink hydrogen:} Electrolysis powered by nuclear energy -- Low Carbon H2.
    \item \textbf{Turquoise hydrogen:} Pyrolysis of natural gas into hydrogen and solid carbon -- Low Carbon H2.
    \item \textbf{White hydrogen:} Naturally occurring geological deposits.
\end{itemize}

Compared to natural gas (CH\textsubscript{4}), hydrogen (H\textsubscript{2}) has a much wider flammability range (4--75\% vs.\ 5--15\%), a lower minimum ignition energy (0.017~mJ vs.\ ~0.3~mJ), a lower volumetric energy density (3.54~kWh/m\textsuperscript{3} vs.\ 11.3~kWh/m\textsuperscript{3}), but a higher gravimetric energy density (39.4~kWh/kg vs.\ 13.3~kWh/kg). Its liquefaction temperature is $-253$\textdegree C, compared to $-162$\textdegree C for LNG, making liquid hydrogen shipping technically challenging and energy-intensive (33\% liquefaction losses vs.\ 10\% for LNG).

\subsection{Belgian Hydrogen Backbone}

Fluxys Hydrogen is developing a phased hydrogen backbone connecting Belgium's industrial valleys (Ghent, Antwerp, Brussels, Charleroi, Liège, Mons) and neighbouring countries (Netherlands, Germany, France, Luxembourg). Key targets include:
\begin{itemize}[noitemsep]
    \item 30~TWh/year transmission capacity by 2030;
    \item First infrastructure operational in 2026 (Zelzate--Kallo and Antwerp South sections, Final Investment Decision already taken);
    \item Cross-border connections with OGE (Germany), Creos (Luxembourg), natTran (France), and HyNetwork Services (Netherlands);
    \item 1,200~km of H\textsubscript{2} pipelines by 2050.
\end{itemize}

At the European level, the European Hydrogen Backbone vision -- developed by 33 infrastructure operators -- targets 31,000~km of H\textsubscript{2} infrastructure by 2030 and 53,000~km by 2040, fed by five supply corridors (North Sea, Nordic and Baltic regions, East and Southeastern Europe, North Africa, and Southwest Europe). The EU27 hydrogen demand is projected at 1,600--2,500~TWh/year by 2050, primarily for industry (700--800~TWh), transport (500--700~TWh), power (200--500~TWh), and buildings (200--500~TWh).

Import of low-carbon hydrogen from outside the EU is expected to be necessary, using carriers such as liquefied hydrogen, ammonia, methanol, LOHC (Liquid Organic Hydrogen Carriers), or synthetic methane. The Belgian Hydrogen Import Coalition (including Fluxys, ENGIE, DEME, EXMAR, and Port of Antwerp-Bruges) has validated the technical and economic feasibility of green hydrogen imports from regions such as Australia, Chile, Oman, Morocco, and Iberia.

Fluxys is also developing an open-access ammonia import terminal (Advario project at Antwerp, targeted operational by 2030), with NH\textsubscript{3} cracking into H\textsubscript{2} and connection to the Fluxys H\textsubscript{2} network. Additionally, the BE-HyStore project, in collaboration with Ghent University, is exploring underground hydrogen storage at Loenhout with a target capacity of 2.4~TWh.

\section{CO\texorpdfstring{\textsubscript{2}}{2} Infrastructure}

\subsection{CO\texorpdfstring{\textsubscript{2}}{2} Transport and the CCUS Value Chain}

Carbon capture and storage/utilisation (CCS/CCU) is increasingly recognised as a necessary mitigation pathway in the EU Industrial Carbon Management Strategy. The EU27 target for CO\textsubscript{2} captured for storage and utilisation is 450~Mt/year by 2050 (from approximately 50~Mt in 2030). CO\textsubscript{2} is transported in dense phase (supercritical, above 74~bara and 31\textdegree C) to maximise pipeline efficiency.

Fluxys's CO\textsubscript{2} value chain covers onshore transmission by pipeline, terminalling (liquefaction, buffer storage, ship loading/unloading, train unloading), and marine shipping to offshore storage or sequestration sites. Sources include industrial carbon capture facilities; destinations include geological storage sites in the North Sea (Norway, UK) and the Netherlands.

\subsection{Belgian CO\texorpdfstring{\textsubscript{2}}{2}Network Development}

The Fluxys c-grid system for Belgium envisions a national CO\textsubscript{2} backbone connecting industrial clusters across Wallonia and Flanders, with ambitions to offer 30~MTPA (million tonnes per annum) of transmission capacity by 2030, with first infrastructure ready by end-2028. The network includes connections towards France (~15--30~Mt/year), Germany and the East (~34--73~Mt/year), and the North Sea (~9--16~Mt/year). Key export projects include:

\begin{itemize}[noitemsep]
    \item \textbf{Zeebrugge offshore CO\textsubscript{2} pipeline} (joint venture with Equinor): 27~MTPA capacity, commissioning target 2031, for sequestration in Norway and/or the UK.
    \item \textbf{Dunkirk Marine CO\textsubscript{2} Export Terminal} (joint venture with Air Liquide and Dunkerque LNG): 1.5~MTPA initial capacity, commissioning target 2028, co-funded by the EU.
    \item \textbf{Antwerp@C Marine CO\textsubscript{2} Export Terminal} (joint venture with Air Liquide and Port of Antwerp-Bruges): up to 10~MTPA, commissioning target 2029, co-funded by the EU.
\end{itemize}

Within the Antwerp cluster, a Final Investment Decision has already been taken for the first pipeline section (Lillo--Kallo--Zandvliet axis), which connects the port area to both Rotterdam (Aramis project in the Netherlands) and Zeebrugge. Construction is partly co-funded by the EU's Connecting Europe Facility (CEF).

\subsection{Tariff Structure and Derisking for CO\texorpdfstring{\textsubscript{2}}{2} Infrastructure}

The development of CO\textsubscript{2} infrastructure faces a chicken-and-egg challenge: the market is nascent, bookings are well below the reference volume (Qref) used in regulated tariff formulas. A high standalone tariff in early years would be unacceptable to first-mover industrial emitters. The proposed solution involves setting an initial tariff aligned with the long-term target tariff, with the resulting revenue gap compensated by government support (e.g., in the form of conditional loans repayable as the market scales up). As capacity bookings grow towards the full Qref, excess revenues cover the repayment. This model mirrors a public-investment logic for new infrastructure with high social value.

\section{The North Sea Integration Model}

Fluxys, in partnership with the University of Liège and funded by Belgium's Energy Transition Fund (SPF Économie), has developed a multi-energy system simulation model covering all 10 North Sea countries (Belgium, France, Germany, Netherlands, Luxembourg, Denmark, Norway, Ireland, the United Kingdom, and Austria). The model optimises on an hourly basis the interactions between electricity, methane (CH\textsubscript{4}), hydrogen (H\textsubscript{2} and derivatives), and CO\textsubscript{2}, targeting carbon neutrality by 2050 at minimum system cost while ensuring energy demand is met at all times.

Key modelling characteristics include linear optimisation over investments and dispatch, perfect foresight (no uncertainty), and no intra-country distance assumption. Inputs include final energy demand (hourly), renewable production potentials (hourly), import assumptions, cost parameters (CAPEX, OPEX, WACC), technical parameters, and CO\textsubscript{2} reduction targets. Outputs include optimal infrastructure investments, hourly dispatch flows, global system cost, marginal prices, energy not served, and CO\textsubscript{2} emissions by country.

Key findings from the North Sea scenario for 2050 include: renewables (solar, wind, hydro) cover electricity demand and green hydrogen production; Norwegian pipeline gas, LNG, and bio-methane cover remaining gas demand and low-carbon hydrogen production via SMR with carbon capture; gas-fired power plants (CCGT and OCGT) run at full capacity only 3\% of the year and supply approximately 1\% of total electricity, but remain indispensable for security of supply during periods of low renewable generation; and CO\textsubscript{2} capture is split equally between biological credits (biomass and bio-CH\textsubscript{4}) and geological capture from low-carbon hydrogen production and industrial processes. An interactive demonstrator of the Belgian model results is publicly accessible at \url{https://integrationdemonstrator.github.io/}.